%% ============================================================
%% main.tex — PNAS research article format
%% Requires pnas-new.cls and pnas-new.bst (both in paper/)
%% Compile: pdflatex main && bibtex main && pdflatex main && pdflatex main
%% ============================================================

\documentclass[9pt,twoside,lineno]{pnas-new}

\templatetype{pnasresearcharticle}

%% ---- Title --------------------------------------------------
\title{Overconfidence in Confidence Interval Estimates Produced by
       Frontier Large Language Models}

%% ---- Authors ------------------------------------------------
\author[a,1]{Author Placeholder}

\affil[a]{Institution Placeholder}

\leadauthor{Walters}

%% ---- Significance statement ---------------------------------
\significancestatement{%
Large language models (LLMs) are increasingly deployed as
decision-support tools, yet their capacity to express calibrated
uncertainty has not been systematically evaluated. We queried three
frontier LLMs across three studies covering equity markets (1 to 22-day
horizons) and six heterogeneous domains. All models were substantially
overconfident, with GPT-4 the most miscalibrated and Claude the least
across all studies. Critically, miscalibration is not uniform: all
models overstate certainty in commodity and equity markets while
overstating uncertainty in cryptocurrency and forex---a domain-confidence
inversion consistent with domain-specific priors inherited from training
corpora. These findings show that global calibration benchmarks are
insufficient; domain-stratified evaluation is necessary to characterize
model-specific calibration risk in high-stakes applications.%
}

%% ---- Author declarations ------------------------------------
\authordeclaration{The authors declare no competing interests.}
\equalauthors{\textsuperscript{1}Correspondence:
  author\texttt{@}institution.edu}
\keywords{large language models | calibration | overconfidence |
  interval estimation | forecasting}

%% ---- Abstract -----------------------------------------------
\begin{abstract}
We document systematic miscalibration in confidence interval estimates
produced by three frontier large language models (Claude, GPT-4, Gemini).
In Study~1, all three models were substantially overconfident when
predicting 1-day equity prices, with GPT-4 exhibiting the largest
Expected Calibration Error (ECE\,=\,43.7\%) and Claude the smallest
(ECE\,=\,19.7\%). In Study~2, using 100 equities across nine prediction
horizons spanning 1--22 days ($N > 18{,}000$ scored predictions), Claude
and GPT-4 became increasingly miscalibrated at longer horizons, whereas
Gemini maintained near-target coverage at the longest horizon tested
(ECE\,=\,2.5\% at 22~days). In Study~3, extending to six heterogeneous
prediction domains, we found a systematic domain-confidence inversion:
all models were overconfident in commodity markets
($\mu = 0.25$ uniformly) but underconfident in cryptocurrency
($\mu = 1.35$--$14.4$), while achieving near-target calibration on
weather. Across all three studies, GPT-4 was the most overconfident
model and Claude the most calibrated. These findings show that LLM
miscalibration is domain- and horizon-structured rather than uniform,
with direct implications for deploying LLMs in high-stakes decision
support.
\end{abstract}

\additionalelement{}
\dates{This manuscript was compiled on \today}
\doi{10.1073/pnas.XXXXXXXXXX}
\acknow{The authors thank the maintainers of the open data APIs used in this
study: yfinance, CoinGecko, Open-Meteo, and ESPN. No external funding
was received for this research.}

%% ============================================================
\begin{document}

\maketitle
\thispagestyle{firststyle}
\ifthenelse{\boolean{shortarticle}}
  {\ifthenelse{\boolean{singlecolumn}}
    {\abscontentformatted}{\abscontent}}{}

%% ---- Body sections -----------------------------------------
% ============================================================
% Introduction
% ============================================================

\section{Introduction}

A well-calibrated forecaster states 80\% confidence intervals that contain
the true outcome 80\% of the time. Decades of research on human judgment
document systematic failure to achieve this standard: people's intervals are
too narrow, a phenomenon termed overconfidence in interval estimates
\citep{Alpert1982, Lichtenstein1982, Klayman1999, Soll2004}. The consequences
are severe. Overconfident interval estimates underlie miscalibrated risk
assessments in medicine \citep{ChristensenSzalanski1981}, finance
\citep{BenDavid2013}, engineering \citep{Flyvbjerg2002}, and strategic
planning \citep{Lovallo2003}.

Large language models (LLMs) are now routinely used as decision-support tools
in precisely these domains, yet whether their probabilistic forecasts are
calibrated remains an open question \citep{Kadavath2022, Lin2022, Xiong2024}.
Prior work has focused on factual question-answering confidence or binary
classification probabilities. Virtually no study has examined LLM calibration
on interval estimates for real-world quantitative forecasting, or asked how
miscalibration varies across prediction horizons and domains.

We address these questions with three preregistered studies. Study~1
establishes a baseline: are three frontier LLMs overconfident when predicting
next-day equity prices with stated confidence intervals? Study~2 extends to
nine prediction horizons (1--22 days) and asks: does miscalibration scale
with forecast horizon? Study~3 extends across six heterogeneous domains
(equities, commodities, cryptocurrency, foreign exchange, weather, NBA game
totals) and asks: is miscalibration domain-specific, and can its direction
reverse?

The human overconfidence literature consistently shows that miscalibration
is domain-specific: experts are most overconfident in their own fields of
ostensible competence \citep{Klayman1999, Soll2004}. It also shows that
overconfidence in interval estimates increases with task difficulty and
time horizon \citep{Griffin1992}. If LLMs inherit domain and horizon
structure from training data, they may replicate or amplify these asymmetries.

Our results confirm and extend the human literature. GPT-4 is the most
overconfident model in every study; Claude is the most calibrated. Study~2
reveals a model divergence at longer horizons: Claude and GPT-4 become
increasingly miscalibrated as the prediction window grows, while Gemini
maintains near-target coverage even at 22~days. Study~3 uncovers a
domain-confidence inversion---all models overconfident in structured financial
markets (equities, commodities) but underconfident in speculative markets
(cryptocurrency, forex)---a pattern absent from prior accounts of LLM
calibration. This paper contributes to the LLM calibration literature
\citep{Kadavath2022, Xiong2024} by focusing on interval estimates rather
than point probabilities and by using real-world outcomes rather than
knowledge-base queries.

% ============================================================
% Study 1 — Baseline overconfidence in 1-day equity predictions
% ============================================================

\subsection*{Study 1}

Study~1 established a baseline calibration estimate for three frontier LLMs
producing confidence intervals for 1-day equity price predictions. We predicted
that all models would be overconfident---stating intervals narrower than warranted
by realized outcomes---and that the degree of miscalibration would differ
systematically across models.

\subsubsection*{Methods}

Three commercially deployed LLMs were evaluated: Claude
(claude-sonnet-4-20250514; Anthropic), GPT-4 (gpt-4o; OpenAI), and Gemini
(gemini-2.5-flash; Google DeepMind). All models were queried via their public
APIs at default temperature with a maximum of 512 output tokens per call.

\paragraph{Stimuli and predictions} Twenty large-cap S\&P~500 equities
were selected as prediction targets (AAPL, NVDA, MSFT, AMZN, GOOGL, META,
AVGO, TSLA, BRK-B, WMT, LLY, JPM, XOM, V, JNJ, MU, MA, ORCL, COST, SPY).
Each prompt provided the current closing price and prediction horizon as plain
text, requesting a point estimate plus 50\%, 80\%, and 90\% confidence
intervals as structured JSON. No chain-of-thought reasoning or calibration
instructions were provided; models were evaluated in their default,
deployment-ready configuration. Predictions were collected on March~23, 2026
for outcomes on March~24, 2026. Each model made 5 independent runs per
ticker, yielding 100 scored predictions per model.

\paragraph{Calibration metrics} For each confidence level
$\alpha \in \{0.50, 0.80, 0.90\}$, the empirical \textit{hit rate}
$\hat{p}_\alpha$ is the proportion of predictions for which the true outcome
fell within the stated interval $[L, U]$. A calibrated model achieves
$\hat{p}_\alpha = \alpha$; values below indicate overconfidence. Expected
Calibration Error (ECE) is the mean absolute deviation between stated and
empirical coverage across all three levels:
\begin{equation}
  \label{eq:ece}
  \text{ECE} = \frac{1}{3} \sum_{\alpha \in \{0.50,\,0.80,\,0.90\}}
               \left| \alpha - \hat{p}_\alpha \right|.
\end{equation}
A perfectly calibrated model achieves ECE\,=\,0.

\subsubsection*{Results}

Table~\ref{tab:study1} summarizes calibration for all three models.

\paragraph{Hit rates} All models were substantially overconfident at every
confidence level. Claude was the least miscalibrated (Hit@80\,=\,58\%,
Hit@90\,=\,73\%), but remained far below calibration targets. GPT-4's 80\%
CI contained the true outcome only 31\% of the time, and its 90\% CI only
38\%---implying interval widths roughly consistent with a 30\% confidence
level. Gemini closely tracked Claude at the 80\% level but underperformed at
90\% (65\%).

\paragraph{ECE} GPT-4's overall ECE (43.7\%) was more than double Claude's
(19.7\%) and more than triple the perfect-calibration target of 0\%.
Gemini was intermediate (22.7\%).

\begin{table}[ht]
  \centering
  \caption{Study~1: Calibration in 1-day equity predictions ($n = 100$ per model).
    ECE = Expected Calibration Error; lower is better.}
  \label{tab:study1}
  \begin{tabular}{lcccc}
    \toprule
    Model   & Hit@50\% & Hit@80\% & Hit@90\% & ECE    \\
    \midrule
    Claude  & 30.0\%   & 58.0\%   & 73.0\%   & 19.7\% \\
    Gemini  & 29.0\%   & 58.0\%   & 65.0\%   & 22.7\% \\
    GPT-4   & 20.0\%   & 31.0\%   & 38.0\%   & 43.7\% \\
    \midrule
    Perfect & 50.0\%   & 80.0\%   & 90.0\%   & 0.0\%  \\
    \bottomrule
  \end{tabular}
\end{table}

\subsubsection*{Discussion}

All three models produced systematically narrow confidence intervals for a
simple, single-domain forecasting task. GPT-4's calibration (31\% coverage at
the 80\% level) is comparable to the most severe human expert overconfidence
on record: \citet{BenDavid2013} found that CFOs' 80\% intervals for annual
equity returns achieved only 36--38\% coverage. These results confirm
universal overconfidence in the simplest possible setting and establish the
model ranking---Claude best, GPT-4 worst, Gemini intermediate---that will
prove consistent across Studies~2 and~3 in the aggregate.

% ============================================================
% Study 2 — Calibration across prediction horizons
% ============================================================

\subsection*{Study 2}

Study~1 established overconfidence at a single, short horizon. Study~2 asked
whether stated confidence intervals---which must widen to maintain coverage
as forecast uncertainty grows---actually scale with prediction horizon.
Using a backtesting design with 100 large-cap equities and nine horizons spanning
1 to 22 trading days, we tested whether miscalibration grows, shrinks, or
diverges across models as the horizon extends.

\subsubsection*{Methods}

The same three models were evaluated using the same prompt format and
structured JSON output as Study~1. Study~2 used a backtesting design in which
predictions were made against historical prices using date-blinded prompts
(current price and horizon provided; calendar date withheld) to prevent
models from conditioning on known future events.

\paragraph{Design} One hundred large-cap S\&P~500 equities were selected.
Predictions were collected at nine horizons: 1, 2, 3, 6, 7, 18, 20, 21, and
22 calendar days. Five independent predictions per ticker per model were
collected at each horizon, using a reference date of March~24, 2026.
Market closures were handled by applying a $\pm$6-calendar-day tolerance
window when matching predictions to realized prices. After excluding
unresolved predictions, the design yielded between 157 and 1,495 scored
predictions per model-horizon cell ($N > 18{,}000$ total).

\paragraph{Calibration metrics} Hit rates and ECE were computed as in
Study~1, aggregated within each model-horizon cell.

\subsubsection*{Results}

Table~\ref{tab:study2} reports 80\% hit rates and ECE at four representative
horizons; complete results for all nine horizons are in the Web Appendix.

\paragraph{Horizon degradation: Claude and GPT-4} Both models exhibited
marked degradation with increasing horizon. Claude's Hit@80 fell from
80.5\% at 1~day to 46.2\% at 18~days (ECE rising from 1.4\% to 29.8\%),
consistent with CI widths anchored to short-horizon volatility that fail
to scale with realized uncertainty at longer windows. GPT-4---already
poorly calibrated at 1~day (ECE\,=\,21.8\%)---deteriorated further, reaching
Hit@80\,=\,27.2\% and ECE\,=\,47.4\% at 18~days. By 22~days, GPT-4's ECE
(49.9\%) approached the theoretical maximum for three-level calibration.

\paragraph{Horizon robustness: Gemini} Gemini exhibited a qualitatively
different pattern. Its 1-day calibration (ECE\,=\,1.8\%, Hit@80\,=\,78.5\%) was
the best of any model at any single horizon in the study. While it degraded
at intermediate horizons (ECE\,=\,16.6\% at 18~days), it recovered strongly
at 22~days (Hit@80\,=\,80.1\%, ECE\,=\,2.5\%)---near-perfect calibration at the
longest horizon tested.

\begin{table}[ht]
  \centering
  \caption{Study~2: Hit@80 and ECE by model and prediction horizon
    (selected horizons; full table in Web Appendix).}
  \label{tab:study2}
  \begin{tabular}{lcccccc}
    \toprule
    & \multicolumn{2}{c}{Claude}
    & \multicolumn{2}{c}{Gemini}
    & \multicolumn{2}{c}{GPT-4} \\
    \cmidrule(lr){2-3}\cmidrule(lr){4-5}\cmidrule(lr){6-7}
    Horizon & H@80   & ECE    & H@80   & ECE    & H@80   & ECE    \\
    \midrule
    1d  & 80.5\% & 1.4\%  & 78.5\% & 1.8\%  & 55.2\% & 21.8\% \\
    6d  & 56.0\% & 20.6\% & 71.7\% & 8.4\%  & 34.5\% & 41.6\% \\
    18d & 46.2\% & 29.8\% & 59.8\% & 16.6\% & 27.2\% & 47.4\% \\
    22d & 51.2\% & 26.9\% & 80.1\% & 2.5\%  & 24.4\% & 49.9\% \\
    \bottomrule
  \end{tabular}
\end{table}

\subsubsection*{Discussion}

Prediction horizon strongly moderated miscalibration in a model-specific
pattern. The degradation of Claude and GPT-4 mirrors the human tendency
toward greater interval underestimation at longer horizons \citep{Griffin1992},
consistent with CI widths anchored to short-run volatility. Gemini's
recovery at 22~days is more surprising: its CI widths appear to scale
approximately proportionally to realized volatility at the extremes of the
tested range but not at intermediate horizons. Whether this reflects a genuine
uncertainty representation or an incidental match with equity return
volatility at these specific windows requires replication at additional
horizons and across different historical periods.

% ============================================================
% Study 3 — Calibration across six prediction domains
% ============================================================

\subsection*{Study 3}

Studies~1 and~2 established overconfidence within a single domain (equities)
and its moderation by horizon. Study~3 asked whether miscalibration is
domain-general or structured by the type of quantity being predicted. Six
heterogeneous domains---spanning financial markets, physical measurements,
and sports outcomes---were included to test whether all models are uniformly
overconfident or whether the direction and magnitude of miscalibration
varies systematically with domain.

\subsubsection*{Methods}

The same prompt format and three confidence levels were used. Note that
Study~3 employed a different Claude model version (claude-sonnet-4-6)
than Studies~1 and~2 (claude-sonnet-4-20250514); GPT-4 (gpt-4o) and
Gemini (gemini-2.5-flash) were identical across all studies (see
Web Appendix for model version details). All predictions used a 24-hour
horizon. Predictions were collected on March~25, 2026; outcomes were
collected automatically on March~26, 2026. One prediction per item per
model was collected.

\paragraph{Domains and stimuli} Six domains were included: equities
($n = 20$ large-cap tickers, same set as Study~1), commodity futures
($n = 5$: gold, crude oil, silver, natural gas, copper), cryptocurrency
($n = 10$ by market cap), foreign exchange ($n = 8$ major pairs), weather
($n = 30$ major U.S.\ cities, next-day high temperature in~$^\circ$F), and
NBA game totals ($n = 3$, subsequently excluded; see Web Appendix). Each
domain prompt provided the current value as context and a question
specifying the target quantity and date. Weather prompts additionally
included the official meteorological forecast for the target day.\footnote{%
  Weather prompts provided the Open-Meteo forecast high temperature for the
  target date. Models produced CIs around (or independent of) this forecast.
  Full prompt templates are in the Web Appendix.}
CoinGecko rate limits caused five cryptocurrency outcomes to remain
unresolved; these were excluded. The final scored dataset comprised
$n = 68$ resolved predictions per model ($N = 204$ total).

\paragraph{Calibration and meta-knowledge metrics} Hit rates and ECE were
computed as in Studies~1 and~2. We additionally computed the Soll--Klayman
meta-knowledge ratio $\mu = \text{MEAD}/\text{MAD}$ \citep{Soll2004},
which quantifies whether stated uncertainty appropriately tracks actual
forecast error. Treating 80\% CI endpoints as the 10th and 90th percentiles
of a symmetric normal ($z = 1.2816$), the Absolute Deviation for prediction
$i$ is:
\begin{equation}
  \label{eq:ad}
  \text{AD}_i = \frac{|(L_i + U_i)/2 - y_i|}{\sigma},
\end{equation}
where $\sigma$ is the within-domain outcome standard deviation. The
meta-knowledge ratio $\mu = \text{MEAD}/\text{MAD}$ is the ratio of
CI-width-implied error (MEAD) to actual midpoint error (MAD) at the
domain level; $\mu < 1$ indicates overconfidence, $\mu = 1$ is perfect,
and $\mu > 1$ indicates underconfidence.

\subsubsection*{Results}

\paragraph{Overall calibration} Aggregating across all 68 resolved items,
Claude was the best calibrated model (ECE\,=\,5.3\%), followed by Gemini
(ECE\,=\,7.2\%) and GPT-4 (ECE\,=\,18.9\%; Table~\ref{tab:study3-overall}).
Overall ECE values are lower than in Studies~1 and~2, driven by
near-perfect calibration on weather---which accounted for 44\% of
Study~3 observations.

\begin{table}[ht]
  \centering
  \caption{Study~3: Overall calibration by model ($n = 68$ resolved
    predictions per model).}
  \label{tab:study3-overall}
  \begin{tabular}{lcccc}
    \toprule
    Model   & Hit@50\% & Hit@80\% & Hit@90\% & ECE    \\
    \midrule
    Claude  & 58.8\%   & 82.4\%   & 85.3\%   & 5.3\%  \\
    Gemini  & 45.6\%   & 73.5\%   & 79.4\%   & 7.2\%  \\
    GPT-4   & 30.9\%   & 63.2\%   & 69.1\%   & 18.9\% \\
    \midrule
    Perfect & 50.0\%   & 80.0\%   & 90.0\%   & 0.0\%  \\
    \bottomrule
  \end{tabular}
\end{table}

\paragraph{Domain effects} Table~\ref{tab:study3-domain} disaggregates
calibration by domain. The central result is a systematic directional
reversal: all models were overconfident in commodity and equity markets but
underconfident in cryptocurrency and, to a lesser degree, forex. No model
was uniformly overconfident or underconfident across all domains.
Commodities showed the most extreme and uniform overconfidence (ECE\,=\,40\%
for all three models; 80\% and 90\% CIs achieving only 40\% coverage).
Crypto showed the opposite pattern: Claude and GPT-4 both achieved 100\%
coverage at the 90\% level. GPT-4's forex calibration was dramatically worse
than the other models (ECE\,=\,44.2\%).

\begin{table}[ht]
  \centering
  \caption{Study~3: Hit rates and ECE by model and domain. $n$ = number
    of resolved predictions. NBA excluded (see Web Appendix).
    Commod.\ = commodity futures.}
  \label{tab:study3-domain}
  \begin{tabular}{llcccccc}
    \toprule
    Domain  & Model  & $n$ & H@50\% & H@80\% & H@90\% & ECE    \\
    \midrule
    Stocks  & Claude & 20  & 50\%  & 85\%  & 90\%  &  1.7\% \\
            & Gemini & 20  & 25\%  & 70\%  & 75\%  & 16.7\% \\
            & GPT-4  & 20  & 15\%  & 55\%  & 65\%  & 28.3\% \\
    \addlinespace
    Commod. & Claude &  5  & 20\%  & 40\%  & 40\%  & 40.0\% \\
            & Gemini &  5  & 20\%  & 40\%  & 40\%  & 40.0\% \\
            & GPT-4  &  5  & 20\%  & 40\%  & 40\%  & 40.0\% \\
    \addlinespace
    Crypto  & Claude &  5  & 80\%  & 100\% & 100\% & 20.0\% \\
            & Gemini &  5  & 20\%  & 80\%  & 100\% & 13.3\% \\
            & GPT-4  &  5  & 20\%  & 100\% & 100\% & 20.0\% \\
    \addlinespace
    Forex   & Claude &  8  & 38\%  &  75\% &  75\% & 10.8\% \\
            & Gemini &  8  & 50\%  &  63\% &  75\% & 10.8\% \\
            & GPT-4  &  8  & 25\%  &  25\% &  38\% & 44.2\% \\
    \addlinespace
    Weather & Claude & 30  & 73\%  &  87\% &  90\% & 10.0\% \\
            & Gemini & 30  & 67\%  &  83\% &  87\% &  7.8\% \\
            & GPT-4  & 30  & 47\%  &  77\% &  80\% &  5.5\% \\
    \bottomrule
  \end{tabular}
\end{table}

\paragraph{Meta-knowledge ratios} Table~\ref{tab:mead} presents the
Soll--Klayman $\mu$ for each model-domain combination. All models show
$\mu = 0.25$ on commodities---stated uncertainty is only one-quarter of
what actual errors require, the most extreme overconfidence observed.
At the other extreme, Claude's crypto $\mu = 14.4$ indicates CIs 14~times
wider than necessary. GPT-4's high forex ECE (44.2\%) is not fully captured
by its $\mu$ (1.29), because $\mu$ cannot distinguish misanchored point
estimates from miscalibrated CI widths; the two metrics must be examined
jointly to diagnose the full failure mode.

\begin{table}[ht]
  \centering
  \caption{Study~3: Meta-knowledge ratio $\mu = \text{MEAD}/\text{MAD}$
    by model and domain. $\mu < 1$ = overconfidence;
    $\mu > 1$ = underconfidence.}
  \label{tab:mead}
  \begin{tabular}{lccc}
    \toprule
    Domain       & Claude & Gemini & GPT-4 \\
    \midrule
    Stocks       & 0.81   & 0.55   & 0.41  \\
    Commodities  & 0.25   & 0.25   & 0.25  \\
    Crypto       & 14.40  & 1.35   & 1.84  \\
    Forex        & 2.68   & 1.67   & 1.29  \\
    Weather      & 1.28   & 1.04   & 0.85  \\
    \midrule
    Overall      & 1.23   & 0.99   & 0.81  \\
    \bottomrule
  \end{tabular}
\end{table}

\subsubsection*{Discussion}

The domain-confidence inversion is the central finding of Study~3: LLMs are
most overconfident in structured financial markets where actual predictability
is low, and most underconfident in speculative markets where they appear to
apply large epistemic-humility priors. This parallels findings from the human
expert literature---\citet{Griffin1992} showed that perceived predictability
often exceeds actual predictability in information-rich, structured
domains---and suggests that LLMs may inherit domain-specific miscalibration
from training corpora dense in equity research and commodity price forecasts.
The uniform $\mu = 0.25$ across all three models on commodities is
particularly striking: this uniformity points to a shared domain prior in
training data rather than a model-specific artifact. The crypto inversion
(Claude $\mu = 14.4$) implies that models apply a ``cryptocurrency is
unpredictable'' narrative that dramatically overstates typical 24-hour
crypto volatility. The Study~3 aggregate model ranking replicates Studies~1
and~2 (Claude best, GPT-4 worst), but aggregate comparisons obscure the
domain-crossing inversions that are the study's primary contribution.

% ============================================================
% General Discussion
% ============================================================

\section{General Discussion}

The most consistent finding across all three studies is that frontier
LLMs---particularly GPT-4---are substantially overconfident in their confidence
intervals for quantitative forecasting. GPT-4's ECE ranged from 18.9\%
(Study~3, favorable domain mix) to 43.7\% (Study~1, equity-only), with a
cross-study mean exceeding 35\%. Even Claude, the best-calibrated model in
every study, achieved only 19.7\% ECE in Study~1. This pattern of overconfidence
was robust to domain, horizon, and sample size, confirming it as a systematic
property of these models rather than an artifact of any single design.

These ECE values are comparable to the most severe human expert overconfidence
on record. \citet{BenDavid2013} found that CFOs' 80\% intervals for annual
equity returns achieved only 36--38\% empirical coverage---roughly equivalent
to our Study~1 Claude result (58\% at 80\%). GPT-4's Study~1 performance
(31\% at the 80\% level) is dramatically worse, suggesting that some frontier
LLMs may exhibit super-human overconfidence in certain forecasting contexts.

\paragraph{The horizon effect} Study~2 reveals that prediction horizon
strongly moderates miscalibration, with a divergent model-specific pattern.
The degradation of Claude and GPT-4 is consistent with CI widths anchored
to short-horizon volatility---a failure mode that mirrors human behavior
at longer horizons \citep{Griffin1992}. Gemini's near-perfect calibration
at both 1~day and 22~days suggests its CI widths scale approximately
proportionally to realized volatility at these extremes, though the mechanism
underlying this unusual profile is not yet understood.

\paragraph{The domain-confidence inversion} Study~3's central finding---that
all models overstate certainty in commodity and equity markets while
overstating uncertainty in cryptocurrency and forex---is inconsistent with a
single domain-agnostic miscalibration mechanism. The most plausible
interpretation is that LLMs apply domain-level heuristics inherited from
training corpora: equity research and commodity analysis provide many narrow
price forecasts, encoding an implicit prior that these markets are predictable,
while cryptocurrency and forex commentary may emphasize unpredictability.
The uniform $\mu = 0.25$ on commodities across all three models supports
a shared training-data origin rather than a model-specific artifact.

\paragraph{Model differences} GPT-4 was consistently the most overconfident
model. Its Study~1 ECE (43.7\%) is nearly double Claude's (19.7\%) in the
same task, and its Study~2 ECE at 22~days (49.9\%) approaches the theoretical
maximum for three-level calibration. Claude performed best in every study at
the aggregate level but exhibits severe directional miscalibration in
cryptocurrency ($\mu = 14.4$)---the largest positive outlier in the
meta-knowledge analysis. This shows that even the best-calibrated model
can be dramatically miscalibrated in specific domains. GPT-4's forex failure
(ECE\,=\,44.2\%) is the worst domain-model result in Study~3 and reflects a
compound failure of point-estimate anchoring and CI-width miscalibration that
the MEAD/MAD ratio alone cannot fully capture; both metrics are required for
complete diagnosis.

\paragraph{Limitations} Several limitations constrain our conclusions.
First, Studies~1 and~3 were each conducted on a single day, and Study~3's
domain-specific results may partly reflect idiosyncratic market conditions
on March~25, 2026. Study~2's backtest design mitigates this concern for
equities but introduces potential look-ahead bias if models' training data
extends to the backtest period. Second, statistical precision is limited in
small Study~3 domains ($n = 5$ for commodities, $n = 5$ for crypto,
$n = 8$ for forex). Third, the Claude model version differed between
Studies~1--2 (claude-sonnet-4-20250514) and Study~3 (claude-sonnet-4-6);
cross-study comparisons for Claude should be interpreted with this caveat.
Fourth, NBA outcomes were excluded entirely from Study~3 due to ESPN API
resolution failures. Fifth, models were evaluated without chain-of-thought
or explicit calibration prompting; whether such interventions produce genuine
calibration improvement is an important open question.

\paragraph{Future directions} Future work should: (i)~collect predictions
daily over extended periods to separate stable calibration properties from
single-day conditions; (ii)~extend the domain set to include earnings
forecasts, macroeconomic indicators, and prediction market prices;
(iii)~evaluate calibration-specific prompting and post-hoc recalibration
(conformal prediction, temperature scaling); and (iv)~use larger $n$ within
Study~3 domains to enable robust domain-level inference.


% -----------------------------------------------------------------------
% Conclusion
% -----------------------------------------------------------------------
\section{Conclusion}

We report three preregistered studies examining whether frontier LLMs exhibit
systematic overconfidence in quantitative interval estimates, and how
miscalibration varies with prediction horizon and domain. All three models
were substantially overconfident across all three studies, with GPT-4
consistently the worst (cross-study ECE 19--44\%) and Claude the best.
Study~2 reveals that Claude and GPT-4 become increasingly miscalibrated as
horizon grows from 1 to 22~days, while Gemini maintains near-target coverage
at both ends of the tested range. Study~3 reveals a domain-confidence
inversion: all models overstate certainty in commodity and equity markets
while overstating uncertainty in cryptocurrency and forex---a pattern
consistent with domain-specific miscalibration inherited from training corpora.

These results have direct practical implications. Users who elicit confidence
intervals from LLMs for financial decision support should not assume that
stated uncertainty reflects calibrated epistemic states. The pattern of
miscalibration is most severe in precisely the domains (commodities, equities)
and at precisely the horizons (multi-week) where overconfident forecasting
carries the most severe consequences. Global calibration benchmarks do not
transfer across domains or horizons; domain-stratified and
horizon-stratified evaluation is necessary to characterize model-specific
calibration risk. More broadly, this study demonstrates that the
domain-specificity and horizon-sensitivity of human overconfidence---developed
over four decades of research on human judges---apply with striking regularity
to large language models trained on those same human outputs.


% -----------------------------------------------------------------------
% Acknowledgements
% -----------------------------------------------------------------------
% Acknowledgements moved to main.tex \begin{acknowledgments} block (PNAS format)


%% ---- Acknowledgements + Bibliography -----------------------
\showacknow
\bibliography{references}

\end{document}
